\chapter{Conclusion and future works}
In this thesis, we presented a multimodal quality inspection framework for the food industry, specifically designed for the bread production process. The framework integrates various techniques, including motion detection, DEIMv2 training, and feature extraction, to assess the quality of dough balls and bread products. The implementation of the framework was carried out using OpenCV, Open3D, Torch, and Tensorboard, demonstrating the practical application of these tools in an industrial setting.
The results obtained from the framework showed promising outcomes in terms of dough-ball volume estimation and bread quality assessment. The quantitative and qualitative results for both Site 1 and Site 2 indicated that the framework is effective in monitoring the production process and identifying potential quality issues.
Future works could focus on further improving the accuracy and efficiency of the framework, as well as exploring additional features and techniques for quality inspection. This could include the integration of more advanced machine learning models, the use of additional sensors for data collection, and the development of real-time monitoring capabilities to provide immediate feedback during the production process. Additionally, expanding the framework to other food products and industries could further demonstrate its versatility and applicability in various manufacturing contexts.